\documentclass[11pt]{article}
\usepackage[a4paper,margin=1in]{geometry}
\usepackage{amsmath,amssymb,amsthm,mathtools}
\usepackage[hidelinks]{hyperref}
\usepackage{booktabs}

\newtheorem{theorem}{Theorem}
\newtheorem{proposition}[theorem]{Proposition}
\newtheorem{lemma}[theorem]{Lemma}
\newcommand{\Tr}{\operatorname{Tr}}
\newcommand{\Q}{\mathbb Q}
\newcommand{\Z}{\mathbb Z}

\title{A Dimension-Four Twisted-Convolution Identity\\
from a Shintani Ray Unit}
\author{Draft for review}
\date{27 July 2026}

\begin{document}
\maketitle

\begin{abstract}
We study the twisted-convolution condition that turns canonical
Shintani--Faddeev real-multiplication values into ghost projectors for
Weyl--Heisenberg symmetric informationally complete measurements.
We give exact Zak-rank, exterior-square, and parity-moment formulations
for the canonical family.  In dimension four, all $36$ shifted ghost
minors have the common factor
$x^2-\sqrt{3+\sqrt5}\,x+1$.  Using the Shintani--Kopp
Kronecker-limit formula and the analytic class-number formula for
$\Q(\sqrt5,\sqrt\phi)$, we identify
$x^2=\phi+\sqrt\phi$, where $\phi=(1+\sqrt5)/2$.
Thus the canonical dimension-four ghost has rank one and satisfies
twisted convolution.  General-dimensional twisted convolution remains
open.
\end{abstract}

\section{Canonical setup}

For $d\ge4$, put
\[
 Q_d=\langle1,1-d,1\rangle,\qquad
 L_d=\begin{pmatrix}d-1&-1\\1&0\end{pmatrix}.
\]
Then
\[
 \operatorname{disc}(Q_d)=(d+1)(d-3),\qquad
 L_d^2+L_d+I=dL_d,
\]
so $L_d^3\equiv I\pmod d$.  Its positive quadratic fixed point is
\[
 \beta_d=\frac{d-1+\sqrt{(d+1)(d-3)}}2.
\]

For $D_{p,q}=\tau^{pq}X^pZ^q$, $\tau=-e^{\pi i/d}$, write
\[
 P=\frac1d\sum_{\boldsymbol p}a_{\boldsymbol p}D_{\boldsymbol p}.
\]
The equation $P^2=P$ is the finite twisted-convolution equation for
the coefficients.  In even dimensions, the representative-wrap signs
of the displacement operators are retained.

\section{Finite rank certificates}

After a scalar shift, shear, and partial Fourier transform, canonical
twisted convolution is equivalent to rank one of a $d\times d$ matrix
$K_d$.  Thus it is equivalent to the vanishing of all $2$-minors.
For any complex matrix $K$,
\[
 \Delta_2(K)=\frac12\left((\Tr K^\dagger K)^2
 -\Tr((K^\dagger K)^2)\right)
 =\sum_{\substack{|I|=2\\|J|=2}}|\det K_{I,J}|^2.
\]
Hence $\Delta_2(K)=0$ if and only if $\operatorname{rank}K\le1$.

The ghost also obeys parity-Hermiticity $G^\dagger=PGP$.  Therefore
$J=PG$ is Hermitian and the rank-one certificate becomes
\[
 \Tr J^4=(\Tr J^2)^2.
\]
Reciprocity forces lower trace information but does not force this
fourth moment.

\section{Dimension four}

Let
\[
 \phi=\frac{1+\sqrt5}{2},\qquad \beta=\phi^2,\qquad
 \tau=-e^{\pi i/4}.
\]
The normalized signed double-sine table is
\[
\begin{pmatrix}
\sqrt5&-x&1&-x^{-1}\\
-x^{-1}&-x^{-1}&-x^{-1}&-x\\
1&-x^{-1}&1&x\\
-x&-x^{-1}&x&-x
\end{pmatrix}.
\]
The zero characteristic is exceptional and is replaced by $\sqrt5$
before applying the displacement chirp.

\begin{proposition}
Every one of the $36$ two-by-two minors of the shifted dimension-four
ghost matrix lies in the principal ideal
\[
 \left(x^2-\sqrt{3+\sqrt5}\,x+1\right)
\]
of $\Q(\sqrt2,\sqrt5)[x,x^{-1}]$.
\end{proposition}

\begin{proof}
The exact Laurent expansion has $34$ formally nonzero minors and two
structurally zero minors.  Reduce each power using
\[
 x^{-1}=t-x,\quad x^2=tx-1,\quad
 t=\sqrt{3+\sqrt5}=\frac{\sqrt2+\sqrt{10}}2.
\]
Every remainder $A+Bx$ is zero.  The deterministic certificate is
generated by
\path{scripts/generate_referee_certificates.py}.
\end{proof}

It remains to prove $x+x^{-1}=\sqrt{3+\sqrt5}$.  Shift and reflection
give
\[
 x=\sqrt2\,
 \frac{S_2(\beta/4\mid\beta,1)S_2(1/4\mid\beta,1)}
 {S_2((\beta+1)/4\mid\beta,1)}.
\]

\section{The modulus-four ray field}

Let $K=\Q(\sqrt5)$ and $\mathcal O_K=\Z[\phi]$.  Since $2$ is inert,
finite residue enumeration gives
\[
 |\operatorname{Cl}_{(4)\infty_2}(K)|=2.
\]
Set $L=K(\sqrt\phi)$.  Its relative discriminant is $(4)$, and the
second real place ramifies because $\phi'<0$.  Hence $L$ is the full
ray field of modulus $(4)\infty_2$.

Put $u=\phi+\sqrt\phi$.  Its Artin conjugate is
$u^\sigma=\phi-\sqrt\phi=u^{-1}$, so
\[
 N_{L/K}(u)=1,\qquad \Tr_{L/K}(u)=1+\sqrt5.
\]

\section{Class number, units, and regulator}

The field $L$ has signature $(2,1)$ and absolute discriminant $-400$.
Its Minkowski bound is $15/(2\pi)<2.4$.  There is no ideal of norm
two, hence $h_L=1$.

Write $t=\sqrt\phi$, so $t^4-t^2-1=0$.  The units $t$ and
$u=t^2+t$ form a fundamental system.  A centered logarithmic
parallelogram argument bounds a hypothetical missing unit at all
archimedean places.  Exact norm enumeration in the integral basis
$1,t,t^2,t^3$ leaves only $\pm1$ in that cell.  Consequently
\[
 \mathcal O_L^\times=\{\pm t^m u^n:m,n\in\Z\},\qquad
 R_L=\log\phi\,\log u.
\]
PARI/GP independently returns discriminant $-400$, class number one,
ray structures $[2]$ and $[2,2]$, and the same regulator; its transcript
is included in \texttt{certificates/pari-audit.txt}.

For the nontrivial quadratic character $\chi$ of $L/K$,
\[
 L(s,\chi)=\frac{\zeta_L(s)}{\zeta_K(s)}.
\]
Since $h_K=h_L=1$, $w_K=w_L=2$, and $R_K=\log\phi$, the analytic
class-number formula at zero yields
\[
 L'(0,\chi)=\log u.
\]

\section{Kronecker-limit normalization}

The two-infinite-place ray group has structure $[2,2]$, while the
one-infinite-place group has structure $[2]$.  The quotient fibers
therefore have order two, so the exponent in Kopp's theorem is $n=1$.
For the two ray classes $A_0,A_1$,
\[
 L(s,\chi)=\zeta(s,A_0)-\zeta(s,A_1).
\]
The generalized derivative in the Shintani--Kopp formula is this
difference, and the normalized positive cocycle occurs squared.
Therefore
\[
 x^2=\exp L'(0,\chi)=u=\phi+\sqrt\phi.
\]

\begin{theorem}
Assuming the published principal-ghost formula and the proved
Shintani--Kopp Kronecker-limit theorem with their stated conventions,
the canonical dimension-four ghost has rank one and satisfies twisted
convolution.  No Stark conjecture is assumed.
\end{theorem}

\begin{proof}
Because $x^2=u$ and $u+u^{-1}=1+\sqrt5$,
\[
 (x+x^{-1})^2=3+\sqrt5.
\]
The defining double-sine value is positive, hence
$x+x^{-1}=\sqrt{3+\sqrt5}$.  The common minor factor vanishes, so all
$36$ minors vanish.  The shifted Zak matrix has rank one, which is
equivalent to normalized ghost idempotency and twisted convolution.
\end{proof}

\section{Scope and review points}

The theorem concerns the canonical dimension-four ghost.  It does not
prove general TCC, the Minimalist RM Values Conjecture, the Stark
conjecture, or Zauner's conjecture in arbitrary dimension.

Before submission, a referee should independently audit:
(i) the reciprocal double-sine convention; (ii) the exceptional zero
characteristic; (iii) the exact $36$-minor certificate; (iv) the ray
class and infinite-place convention; (v) the cocycle-square
normalization in Kopp's theorem; and (vi) the final projector
normalization.

\begin{thebibliography}{9}
\bibitem{Kopp}
G.~Kopp, \emph{The Shintani--Faddeev modular cocycle: Stark units from
$q$-Pochhammer ratios}, arXiv:2411.06763.
\bibitem{Shintani1977}
T.~Shintani, \emph{On a Kronecker limit formula for real quadratic
fields}, J. Fac. Sci. Univ. Tokyo 24 (1977), 167--199.
\bibitem{Shintani1978}
T.~Shintani, \emph{On certain ray class invariants of real quadratic
fields}, J. Math. Soc. Japan 30 (1978), 139--167.
\bibitem{ZaunerJL}
S.~Flammia et al., \emph{Zauner.jl},
\url{https://github.com/sflammia/Zauner.jl}.
\end{thebibliography}

\end{document}
